% Round 15. Scope analysis and a formal consequence of frozen R13--R14.
% The formal HB1--HB7 proof is unchanged; the scope discussion is updated
% to the completed arbitrary-base and all-rank results of this round.
\section{Degree-two center normalization and the fourfold boundary}
\label{r15:higher-center-boundary}

This section records the formal center-normalization step first
isolated during the relative curve calculation.  Its final scope
discussion incorporates the subsequent arbitrary-base and all-rank
results of this round. The connecting sequence used below is the frozen Round~13
result, not a new construction. Combining it with the subsequently
proved Round~14 constant-fiber-dimension injectivity gives an additional
formal corollary. The Round~13 source states the kernel quotient only
under $r>\dim C$; the searched Round~14 sources retain that dimension
condition when quoting the quotient. The argument below records exactly
why it can now be removed in degree $r=2$, without treating the raw
exceptional restriction as closed.

Let $X/K$ be smooth integral projective, let $j:C\hookrightarrow X$
be a nonempty smooth integral closed center of codimension $c\geq2$,
and write
\[
 q:Y=\operatorname{Bl}_C X\longrightarrow X,\qquad
 i:E=\mathbf P_C(N_{C/X})\hookrightarrow Y,\qquad
 \pi:E\longrightarrow C.
\tag{HB1}
\]
The convention for projectivization is that of lines, so
$\mathcal O_Y(E)|_E=\mathcal O_E(-1)$.
A disconnected smooth center is handled componentwise. All forms,
finite maps and degrees are native over the original $K$.

\begin{lemma}[A primitive on the center is already available]
\label{r15:center-primitive}
Suppose that $\Omega\in B^{2,2}(X)$ is separately closed and that
$\beta\in B^{1,1}(Y)$ satisfies $D\beta=q^*\Omega$.
Then there is $\zeta\in B^{1,1}(C)$ with
\[
                         D\zeta=j^*\Omega.
\tag{HB2}
\]
One explicit choice is
\[
 \zeta=\frac1e\operatorname{Tr}_f(h^*i^*\beta),\qquad
 h:T\longrightarrow E,\quad f=\pi h:T\longrightarrow C
 \text{ finite dominant},\quad e=\deg f,
\tag{HB3}
\]
where $T$ is a finite normal multisection of the projective bundle.
No actual proper trace along $\pi$ is used.
\end{lemma}

\begin{proof}
The map $\pi$ is projective, smooth, surjective, and all its geometric
fibers have dimension $c-1$. Source and target are normal integral
projective varieties. Thus the frozen Round~14 theorem
\texttt{foundations\_audit/equidimensional\_projective\_descent.tex},
label \texttt{r14j:equidimensional} (assembled Theorem~16.1),
applies. In particular $\pi^*$ is injective on actual
$H_{\mathrm{BC}}^{2,2}(C)$. Restriction of the supplied primitive gives
\[
                  \pi^*j^*\Omega=D(i^*\beta).
\]
This already proves existence in (HB2).

For the explicit formula, the finite-normal-multisection lemma
\texttt{r14e:finite-multisection} gives $h,T,f$ as in (HB3), and
actual pullback gives
\[
              f^*j^*\Omega=D(h^*i^*\beta).
\]
The form inside the trace has bidegree $(1,1)$, within the canonical
actual finite strip trace. Its derivative-current identity and the
native finite degree formula give
$D\operatorname{Tr}_f(h^*i^*\beta)=e\,j^*\Omega$ as actual forms.
This is exactly Corollary \texttt{r14j:primitive}, formula
\texttt{r14j:primitive-formula} (assembled Corollary~16.2),
applied to $\pi$. All tests in the finite trace remain actual tests
on the original $C$.
\end{proof}

\begin{definition}[The normalized exceptional class]
\label{r15:normalized-exceptional-class}
For the supplied pair $(\Omega,\beta)$ and a choice of (HB2), put
\[
 \theta_{\beta,\zeta}=i^*\beta-\pi^*\zeta\in B^{1,1}(E),
 \qquad
 \mathcal Q_C=
 \frac{H_{\mathrm A,\mathrm{act}}^{1,1}(E)}
 {i^*H_{\mathrm A,\mathrm{act}}^{1,1}(Y)
       +\pi^*H_{\mathrm A,\mathrm{act}}^{1,1}(C)}.
\tag{HB4}
\]
Then $D\theta_{\beta,\zeta}=0$.
Its image $\varepsilon_q([\Omega])$ in $\mathcal Q_C$ is the
normalized exceptional class. Defining $\theta$ requires no extension
of $\zeta$ off the center.
\end{definition}

\begin{proposition}[The degree-two exceptional quotient for every smooth center]
\label{r15:degree-two-all-centers}
In (HB1), the frozen connecting map induces a canonical isomorphism
\[
 \mathcal Q_C\ \xrightarrow{\ \sim\ }\
 \ker\bigl(q^*:H_{\mathrm{BC,act}}^{2,2}(X)
                  \longrightarrow H_{\mathrm{BC,act}}^{2,2}(Y)\bigr).
\tag{HB5}
\]
The class of (HB4) is independent of the supplied primitive and of the
center primitive, and corresponds to $[\Omega]$ under (HB5).
In particular, for the supplied $\Omega$, it vanishes if and only if
$\Omega$ admits a global actual $D$-primitive on $X$.
This is a statement about primitive descent, not a converse to the
sufficient implication $\mathfrak o_X^2([Z])=0\Rightarrow G(Z)$.
\end{proposition}

\begin{proof}
The frozen source
\texttt{round13/analytic\_descent/retracted\_modification\_descent.tex},
Theorem \texttt{r13d:connecting}, equation (R13.D16), gives exactness
at the middle two terms of
\[
 H_{\mathrm A}^{1,1}(Y)\oplus H_{\mathrm A}^{1,1}(C)
 \xrightarrow{i^*-\pi^*}H_{\mathrm A}^{1,1}(E)
 \xrightarrow{\partial_D}H_{\mathrm{BC}}^{2,2}(X)
 \xrightarrow{(q^*,j^*)}
 H_{\mathrm{BC}}^{2,2}(Y)\oplus H_{\mathrm{BC}}^{2,2}(C).
\tag{HB6}
\]
It holds without the condition $2>\dim C$. Exactness at the Aeppli
term says that the full kernel of $\partial_D$ is precisely the sum
of the two images in the denominator of (HB4); the minus sign does
not change their sum as vector subspaces.

If $q^*[\Omega]=0$, the commutative square gives
$\pi^*j^*[\Omega]=i^*q^*[\Omega]=0$. Lemma~\ref{r15:center-primitive},
or the injectivity used in its proof, gives $j^*[\Omega]=0$.
Consequently $\ker q^*=\ker(q^*,j^*)$. The image statement in (HB6)
therefore proves (HB5).

For the representative check, extend $\zeta$ degreewise to an actual
$\eta\in B^{1,1}(X)$ using the frozen immersion-germ extension theorem.
Put
\[
 \Omega_0=\Omega-D\eta,\qquad
 \beta_0=\beta-q^*\eta.
\tag{HB7}
\]
Then $j^*\Omega_0=0$, $D\beta_0=q^*\Omega_0$, and
$i^*\beta_0=\theta_{\beta,\zeta}$. The connecting-map construction
with the extension $\beta_0$ therefore sends $[\theta]$ to
$[\Omega_0]=[\Omega]$. This explains precisely where extending
$\zeta$ changes the curvature representative and why its
Bott--Chern class is unchanged.

Replacing $\beta$ changes $\theta$ by the restriction of a global
$D$-closed form on $Y$. Replacing $\zeta$ changes it by the pullback
of a $D$-closed form on $C$. Both changes lie in the complete
quotient denominator. If $\Omega$ is changed by $D\xi$, take
$\beta+q^*\xi$ and $\zeta+j^*\xi$; the two added restrictions cancel
in $\theta$. Thus the normalized class is canonical. The last
assertion follows either from (HB5) or from the explicit corrected
primitive formula in the frozen connecting proof.
\end{proof}

\subsection{The three fourfold exceptional cases are now controlled}

Now let $\dim X=4$ and retain curvature degree $(2,2)$.
A nontrivial smooth-center blowup has center dimension $0$, $1$, or
$2$; a smooth Cartier-center blowup is an isomorphism.
Let
\[
 h_E=c_1(\mathcal O_E(1))=-i^*c_1(\mathcal O_Y(E))
 \quad\text{in }H_{\mathrm A}^{1,1}(E),
\tag{HB8}
\]
with compatible model metrics, or with their scalar $D$-differences
absorbed in the Aeppli class. Hence the normal hyperplane class is
already in the first denominator of (HB4).

\begin{enumerate}
\item \emph{Point center.} Here $E\simeq\mathbf P^3_{\kappa(p)}$.
The raw restriction $i^*\beta$ is $D$-closed and the center contributes
no positive-degree form. The completed all-rank projective-bundle
Aeppli theorem gives
$H_{\mathrm A,act}^{1,1}(\mathbf P^3)=\mathbf R[h_E]$,
and this generator is already a restriction from $Y$.
This conclusion uses the full incidence argument of
Section~\ref{r15i:section}; it is not inferred from the
$\mathbf P^2$ computation by restriction to a hyperplane.

\item \emph{Curve center.} Here $E=\mathbf P_C(N_{C/X})$ is a
$\mathbf P^2$-bundle over a smooth curve. Again $j^*\Omega=0$ by
dimension, so $\zeta=0$ is available. The completed relative theorem is
\[
 H_{\mathrm A}^{1,1}(E)
       =\pi^*H_{\mathrm A}^{1,1}(C)+\mathbf R[h_E].
\tag{HB9}
\]
The center summand is one-dimensional by the frozen proper-curve
potential theorem. The proof of (HB9) uses the ordinary PWS
projective-bundle integration theorem, all smooth mixed-test
pairings, algebraic base change, and the global incidence
homotopy. It does not infer a family statement merely from
the groups of individual fibers. Both summands are now in
the denominator of (HB4).

\item \emph{Surface center.} Here $E$ is a $\mathbf P^1$-bundle over
the smooth projective surface $C$. In general
$D(i^*\beta)=\pi^*j^*\Omega$ is nonzero. One must first construct
$\zeta$ by (HB2)--(HB3) and use $\theta=i^*\beta-\pi^*\zeta$.
That center-primitive step is already proved; it is not an additional
open existence conjecture. The now proved relative formula
(HB9), with the full possibly unknown group
$H_{\mathrm A}^{1,1}(C)$, kills (HB4). The necessary integration
and Green outputs hold on arbitrary-dimensional smooth bases,
and the bundle statement glues over its algebraic
trivializations. Computing $H_{\mathrm A}^{1,1}(C)$ itself
is not required by this relative formulation.
\end{enumerate}

Thus none of these three minimum-one exceptional Aeppli
problems remains open in the completed round.  In fact
Theorem~\ref{r15m:strip-injectivity} proves smooth-center
Bott--Chern injectivity whenever the smaller bidegree is
at most two, in every dimension, using the full first-strip
formula and localized positive normal Chern extensions.
This is a primitive-descent theorem.  It does not on its
own supply a numerically admissible model witness for
every codimension-two cycle on a fourfold.  The remaining
model question is the primitive numerical residual in
Section~\ref{r15res:section}; no universal vanishing of
that residual or general strong-G conclusion is claimed.

\subsection{The historical fixed-tuple obstruction and its local resolution}

The first fixed-global-tuple proof had the following genuine
limitation. Over a curve, the normalization in a finite splitting field is again
smooth in characteristic zero. The poles of the root sections, their
pairwise collision loci, and the vanishing loci of finitely many
coefficient functions form a finite set of closed points. After a
finite constant extension those points are rigid rational points.
The frozen immersion-germ theorem and properness give
$\alpha|_{U_b\times\mathbf P^1}=\operatorname{pr}_z^*\beta_b$ on a
whole ordinary product neighborhood. The one-variable convolution
then supplies all mixed test pairings there.

For a base of dimension at least two, normalization in that same
splitting field is generally only normal, not smooth. The bad locus
can contain divisors and their intersections, and rational root maps
to $\mathbf P^1$ can have higher-codimension indeterminacy. Thus it
is not a finite collection of rigid fibers. Restricting an input to
$D\times\mathbf P^1$ along a bad divisor produces a relative form
with a positive-dimensional base; it is not a fixed one-variable
form $\beta_b$. The immersion-germ extension theorem gives
extensions of form germs, but no analytic retraction from a whole
ordinary neighborhood of $D$ to $D$, and no prescribed-current
compatibility for an arbitrary extension. It cannot substitute for
the missing family operator.

Resolving that fixed normal splitting cover or its root-map
indeterminacy would introduce proper modifications. The finite
PWS trace alone would not descend an output through their
positive-dimensional fibers. These observations remain correct
as a diagnosis of the earlier global-tuple strategy.

The subsequent local separable-moment perturbation argument
does not make that resolution step. Around each arbitrary
analytic base point it first takes finite complete-fiber
Laurent data, then perturbs the finitely many polynomial
coefficients by sufficiently small elements of $K$.
On the compact domains the norms of every relevant unit
and every domain inequality are unchanged exactly, while
the specialized leading coefficients, discriminants, and
pairwise resultants are nonzero. Consequently the splitting
cover is finite etale on an ordinary neighborhood of the
chosen point. There is no remaining bad divisor on that
local neighborhood. Ordinary gluing follows from uniqueness
of the prescribed currents, not from an assumed retraction
from a neighborhood of a bad divisor.

This is the local input used in
Theorem~\ref{r15hp:projection} and
Proposition~\ref{r15hp:base-change}; its projective-bundle
extension is Theorem~\ref{r15ps:projective-bundle}.
Therefore the fixed-global-tuple limitations above are
historical method boundaries, not unresolved hypotheses
of the completed arbitrary-base PWS output theorem.
The scope still does not include arbitrary proper
pushforward of positive-delta-codimension actual forms.
